\section{实验方法}
\label{sec:method}

\subsection{硬件与工具链}

实验平台为 AMD/Xilinx KV260，器件 XCK26-SFVC784-2LV-C。RTL 使用
Verilog，Vivado/Vitis 版本均为 2022.2；正式 profile 为
\code{abi\_v2\_release\_200}，目标时钟 200 MHz，实际 PL 时钟
199.998 MHz。阵列参数为 ROWS=18、COLS=16、COUT\_TILE=32，AXI-Stream
宽度 64 bit。正式 bitstream 由 clean Git commit 和 fresh build 生成，
post-place 与 post-route 均执行时钟、资源、拥塞、未约束路径、路由完整性、
DRC、setup/hold/pulse 和内部 endpoint 门禁。报告中的功耗若无特殊说明均不
表示板级测量。本文使用正式 routed DCP 的 Vivado post-route estimated power：
process 取 typical，环境温度 25\,$^\circ$C，vectorless toggle rate 为 12.5\%，
static probability 为 0.5，未加载 SAIF；报告同时给出结温和活动率置信度。

软件运行于 Cortex-A53 裸机 EL1，使用四核协同和 lwIP 持久 TCP runner。
COCO 输入为 batch=1、416$\times$416 固定方形 letterbox：RGB、保持比例、
PIL bilinear、fill=114、居中 padding。FP32、INT8 host 和板端读取完全相同的
预处理 uint8 字节与 metadata。accuracy decode 固定 conf=0.001、NMS IoU=0.65、
multi-label、class-aware、max\_nms=30000 和 max\_det=300。

\subsection{模型与数值基线}

上游资产为 Ultralytics v9.5.0 YOLOv3-tiny COCO80，官方权重文件 SHA256
为 \texttt{74fb61c9...\allowbreak 946280}。FP32-640 按上游代码和 COCO val2017 复现，用于
确认资产；FP32-deploy416 使用本文固定预处理和后处理，作为唯一量化损失
母线。当前部署 INT8 采用 Conv--BN 融合、per-tensor symmetric int8 weight
（zero point=0）和 per-tensor affine reduced-range activation；13 个卷积均
导出 multiplier/shift/zero-point、int32 bias 和 256 项 activation LUT。

精度验证分为四级：PyTorch INT8 检查量化方案；RTL-semantic host 以最终
packed 参数精确模拟中心化、int8 MAC、int32 bias、RTL rounding/clamp/LUT
和 A53 tensor op；板端 conformance 与其逐字节比较；COCOeval 对 raw head
统一执行 canonical host decode/NMS。只有板端 raw head 与 host 完全相同，
才能把 network mAP 等同于 host INT8 mAP。product mAP 则由 A53 实际
detections 生成，单独度量后处理实现差异。

\subsection{正确性样本与故障注入}

逐节点 conformance 使用 128 张代表图，覆盖 80 类、空图、拥挤场景、小/
中/大目标和不同宽高比。每张保存输入、13 个卷积、6 个池化/分支节点、
upsample、concat 和双 head，共 22 个张量。比较项目为 shape、byte count、
完整 payload SHA 和逐字节 mismatch；任一字节不一致均失败。另使用 synthetic
测试覆盖全零点、min/max、累加边界、m19 尾 tile、255 尾通道、TKEEP、
无候选、多标签同框和 NMS 阈值相等。

网络通道先以 1、128 和 5000 张纯传输验证 CRC、sequence、image id、顺序、
DDR 地址、cache 和 DMA 可见性；再注入错误长度、CRC、重复/跳跃序号、错误
参数 hash 和断线。板端必须发送 ERROR 并终止 session，不得继续消费后续
chunk。稳定性测试循环固定 128 张集合 600 s，要求输出 CRC 稳定、协议/DMA/
PL/timeout 错误为零；能读取温度时峰值严格低于 85\,$^\circ$C。

\subsection{性能口径}

主性能独立运行 3 次，每次 20 个 warmup 和 1000 个 timed sample。参数常驻
DDR，关闭 raw dump、逐图 UART 和 JSON 格式化。timer 由约 100 MHz 的通用
计数器驱动，并与 200 MHz PL cycle counter 交叉检查。记录项目包括 13 个
PL dispatch、PL 总时间、普通 pool、特殊 pool、requant-concat、P5 copy、
candidate、sort、NMS、decode、resident 和 chunk pipeline。

报告 min、mean、P50、P90、P95、P99、max 和 FPS。resident 从输入量化张量
驻 DDR 到 detections 完成；network resident 还可以在 raw head 处停止；
postprocess 从双 raw head 到 detections；pipeline 额外包含 TCP 分块发送和
结果返回。SD I/O 和 cold boot 另表，不能与 resident 相加后仍称“加速器
延迟”。当前第一版不以 FPS 作为模型 release 门禁，但所有 3000 条记录必须
完整、输出确定、分项与总时长自洽。

统计时每次独立运行重新建立网络会话并检查参数与输出身份，但不重新生成
bitstream 或参数，以隔离部署启动与稳态计算。warmup 用于稳定 cache、分支
预测和协议队列，不进入分位数；三次运行分别保留原始逐图 CSV，汇总脚本从
样本而非预先计算的均值重算 P50/P90/P95/P99。FPS 使用总 timed wall time
和样本数计算，同时报告 $1000/\overline{T}$ 作为一致性检查。若两者因丢样、
重试或时间边界不同而不一致，运行判为失败而不是选择较好数字。

单张 resident 计时只覆盖输入已在 DDR 后的计算边界，适合比较软硬件调度；
pipeline 计时覆盖主机分块、TCP、板端计算和结果返回，适合评估验证工具吞吐；
SD I/O 则受文件系统和卡速影响。三种口径回答不同问题，不能相加或互换。
尤其是 UART 输出必须位于 timer 之后：115200 baud 下数百字节文本即可达到
毫秒级，足以掩盖真实 decode 计算。本文因此既保存分项 cycle，也保存主机
chunk wall time，以检查局部计时没有遗漏系统瓶颈。

性能有效性还要求温度与错误状态稳定。每个 timed record 绑定输出 CRC、PL
错误计数、DMA 状态和 timeout；任一异常会使整次运行失效，不能只删除该点。
600 s soak 检查长期重放和网络会话稳定，但它不替代 3000 样本的分布统计；
冷启动复测验证 BOOT 路径不改变内核结果，也不替代正式网络/JTAG主结果。

\subsection{CPU 基线与受控消融}

CPU 基线必须在同一 KV260、同一 416 输入、同一量化参数和同一 A53 后处理
下运行完整 13 卷积 DAG。单核标量 C 版本不使用 NEON；四核版本以输出通道
块并行并使用 NEON，参数常驻 DDR，排除 TCP/SD/UART。二者与 \LASA 使用同一
RTL-semantic rounding/LUT，输出需逐节点 byte-exact。报告 resident、卷积、
tensor op 和 decode，以及 3$\times$1000 的相同统计。

受控消融要求所有 bitstream 来自同一最终 RTL 根、同一 18$\times$16 阵列、
同一 200 MHz 约束和 clean provenance。变体包括：原生 1$\times$1 对稀疏
3$\times$3；逐 pass IFM 对 HWC materialize/replay；无预载/无 handoff、仅
预载、预载+handoff、完整 overlap；保守固定 tile 对层自适应 tile；单核对
四核 A53。每个变体先通过 byte-exact，再报告 latency、DDR bytes、compute
idle、context gap 和资源。不能执行完整网络的变体只报告代表层，且不得
外推端到端加速。

\todo{完成同板单核标量 C 与四核 NEON/output-channel 基线。当前仓库只有
基线合同和结果表生成位置，没有可以引用的最终数值。}

\todo{完成所有受控 bitstream 消融，并把每个 DCP/bit、Git SHA、门禁报告、
计数器 CSV 和 byte-exact 报告纳入统一 manifest。}

\subsection{外部工作对比原则}

外部表格固定列出器件、频率、网络/输入、batch、是否完整双尺度、是否包含
后处理、延迟、FPS、精度和资源。由于不同论文对“GOPS”“latency”和功耗的
定义差异很大，本文不把 conv-only GOPS 与端到端 FPS 直接排序，也不把
Vivado 估算功耗与板级测量混算 GOPS/W。比较重点是功能完整性、层适应性、
验证深度和真实板级 latency。引用数值必须回到原始论文，不从综述二次抄录。

实验设计还区分内部效度、构念效度与外部效度。内部效度由同源 RTL、clean
provenance、相同输入和 byte-exact 控制，防止加速来自功能删减；构念效度由
resident、pipeline、GOPS 与 FPS 的明确边界控制，防止不同指标被混称；外部
效度则受单一器件、固定416输入、per-tensor INT8 与 YOLOv3-tiny 网络限制。
本文不会把一个 XCK26 上的结果外推为所有 FPGA 或所有 CNN，也不会用当前
量化参数的精度代表架构上限。公开 artifact 的目标是让读者复算本文数据，
而不是消除跨平台比较本身的口径差异。
